\section{System design}
\label{sec:system}

FallGuard performs pose extraction and temporal classification on the Raspberry Pi. The Hailo accelerator runs the pose network, while the Pi CPU evaluates a one-second skeleton window with MaskedBiMamba. The local alert state and prediction records remain available during network outages. When connectivity returns, a synchronization service transfers pending records to the cloud for event history and dashboard updates. Figure~\ref{fig:edge_flow} shows this data flow.

\begin{figure*}[tp]
\centering
\begin{tikzpicture}[font=\small,box/.style={draw,rounded corners=2pt,align=center,minimum height=0.85cm,text width=2.25cm,fill=blue!4},arr/.style={-{Latex[length=2mm]},thick}]
\node[box] (camera) {Camera or local video};
\node[box,right=0.35cm of camera] (pose) {YOLOv8 pose\\Hailo / CPU};
\node[box,right=0.35cm of pose] (model) {1-s pose window\\MaskedBiMamba};
\node[box,right=0.35cm of model] (queue) {Local alert state\\SQLite records};
\draw[arr] (camera)--(pose);\draw[arr] (pose)--(model);\draw[arr] (model)--(queue);
\node[box,below=1.2cm of queue,fill=orange!6] (cloud) {Cloud event store\\and dashboard};
\draw[arr] ([xshift=-.35cm]queue.south)--node[left,font=\footnotesize,align=right]{Pending records\\on reconnection}([xshift=-.35cm]cloud.north);
\draw[arr] ([xshift=.35cm]cloud.north)--node[right,font=\footnotesize]{Acknowledgement}([xshift=.35cm]queue.south);
\begin{scope}[on background layer]
\node[draw,dashed,rounded corners,fit=(camera)(pose)(model)(queue),inner sep=8pt,label=above:{Raspberry Pi: inference continues without a network}] {};
\end{scope}
\end{tikzpicture}
\caption{Local inference and deferred cloud synchronization. CPU pose estimation provides a fallback when the Hailo backend is unavailable. Cloud acknowledgement follows durable storage, allowing safe retries after an interrupted transfer.}
\label{fig:edge_flow}
\end{figure*}

Routine synchronization transfers pose features and classification records. An authenticated preview can relay annotated images for monitoring; preview frames are held temporarily and kept separate from the event archive.

\subsection{Alert-state transitions}

The edge node converts a sequence of fall probabilities into \textsc{Normal}, \textsc{Suspected}, \textsc{Confirmed}, and \textsc{Cleared} states. Let $p_k$ be the probability of the $k$th locally classified window. The high- and low-probability counters are
\begin{equation}
\begin{aligned}
 c_k^+ &= \mathbf{1}[p_k\geq0.65](c_{k-1}^++1),\\
 c_k^- &= \mathbf{1}[p_k\leq0.35](c_{k-1}^-+1).
\end{aligned}
\label{eq:alert_counters}
\end{equation}
A high observation opens a suspected event. Confirmation requires $c_k^+\geq2$ and $p_k\geq0.82$; $c_k^-\geq2$ clears an active event. Requiring consecutive high-probability observations prevents an isolated spike from confirming an event. Acknowledgement and user notes are stored with the event, separately from the prediction sequence.
